\documentclass[11pt,a4paper]{article}

% ── Packages ──────────────────────────────────────────────────────────────────
\usepackage[utf8]{inputenc}
\usepackage[T1]{fontenc}
\usepackage[margin=2.5cm]{geometry}
\usepackage{hyperref}
\usepackage{xcolor}
\usepackage{listings}
\usepackage{enumitem}
\usepackage{booktabs}
\usepackage{longtable}
\usepackage{graphicx}
\usepackage{amsmath,amssymb}
\usepackage{fancyhdr}
\usepackage{titlesec}

% ── Colors ────────────────────────────────────────────────────────────────────
\definecolor{isls-blue}{HTML}{1a5276}
\definecolor{isls-orange}{HTML}{e67e22}
\definecolor{isls-green}{HTML}{27ae60}
\definecolor{isls-red}{HTML}{c0392b}
\definecolor{isls-gray}{HTML}{7f8c8d}
\definecolor{code-bg}{HTML}{f7f9fb}

% ── Listings ──────────────────────────────────────────────────────────────────
\lstdefinestyle{rust}{
  language=C, % closest built-in
  basicstyle=\ttfamily\small,
  keywordstyle=\color{isls-blue}\bfseries,
  commentstyle=\color{isls-gray}\itshape,
  stringstyle=\color{isls-green},
  backgroundcolor=\color{code-bg},
  frame=single,
  rulecolor=\color{isls-gray!30},
  breaklines=true,
  showstringspaces=false,
  tabsize=4,
  morekeywords={fn,let,mut,pub,struct,impl,enum,use,mod,self,Self,
    crate,super,where,trait,for,in,if,else,match,return,async,await,
    Ok,Err,Some,None,Vec,String,HashMap,HashSet,PathBuf,Result,Option},
}
\lstdefinestyle{bash}{
  language=bash,
  basicstyle=\ttfamily\small,
  backgroundcolor=\color{code-bg},
  frame=single,
  rulecolor=\color{isls-gray!30},
  breaklines=true,
  showstringspaces=false,
}
\lstdefinestyle{toml}{
  basicstyle=\ttfamily\small,
  backgroundcolor=\color{code-bg},
  frame=single,
  rulecolor=\color{isls-gray!30},
  breaklines=true,
  showstringspaces=false,
  commentstyle=\color{isls-gray}\itshape,
  morecomment=[l]{\#},
}

% ── Headers ───────────────────────────────────────────────────────────────────
\pagestyle{fancy}
\fancyhf{}
\fancyhead[L]{\small\textcolor{isls-gray}{ISLS D4 Spec --- Ophan-Orbits}}
\fancyhead[R]{\small\textcolor{isls-gray}{v1.0.0 --- \today}}
\fancyfoot[C]{\thepage}

% ── Section styling ───────────────────────────────────────────────────────────
\titleformat{\section}
  {\Large\bfseries\color{isls-blue}}{\thesection}{1em}{}
\titleformat{\subsection}
  {\large\bfseries\color{isls-blue!80}}{\thesubsection}{1em}{}
\titleformat{\subsubsection}
  {\normalsize\bfseries\color{isls-blue!60}}{\thesubsubsection}{1em}{}

\hypersetup{
  colorlinks=true,
  linkcolor=isls-blue,
  urlcolor=isls-orange,
  citecolor=isls-green,
}

% ══════════════════════════════════════════════════════════════════════════════
\begin{document}

% ── Titlepage ─────────────────────────────────────────────────────────────────
\begin{titlepage}
\centering
\vspace*{3cm}

{\Huge\bfseries\color{isls-blue} ISLS D4 Specification}\\[0.5cm]
{\LARGE\color{isls-orange} Ophan-Orbits: Auto-Norm Emergence}\\[2cm]

{\large
\begin{tabular}{rl}
\textbf{System:}    & ISLS --- Intelligent Semantic Ledger Substrate \\
\textbf{Version:}   & v3.5 (D4) \\
\textbf{Author:}    & Sebastian Klemm \\
\textbf{Date:}      & \today \\
\textbf{Status:}    & Specification \\
\textbf{Depends:}   & D3 (complete), isls-norms v3.0 \\
\textbf{Branch:}    & \texttt{claude/implement-isls-d4-spec} \\
\end{tabular}
}

\vfill

{\small\color{isls-gray}
Mithras Substructure University\\
Repository: \texttt{https://github.com/LashSesh/graphity}\\
License: MIT
}
\end{titlepage}

% ── TOC ───────────────────────────────────────────────────────────────────────
\tableofcontents
\newpage

% ══════════════════════════════════════════════════════════════════════════════
\section{Executive Summary}
% ══════════════════════════════════════════════════════════════════════════════

D4 closes the feedback loop in the ISLS generation pipeline. After D3, ISLS 
generates applications from natural language, but every generation starts from 
zero --- the system learns nothing from prior runs. D4 introduces 
\textbf{auto-norm emergence}: the system observes what it generates, detects 
recurring cross-layer patterns across domains, and promotes stable patterns to 
first-class norms that influence future generations.

In Dual Fabric / TriM\"obius terminology: each synthesis run produces an 
\emph{Ophan orbit} --- a coherent rotation through the full-stack layers 
(Database $\to$ Model $\to$ Query $\to$ Service $\to$ API $\to$ Frontend 
$\to$ Test). When multiple orbits across different domains trace the same 
topology, the shared structure crystallises into a norm. The norm is the 
\emph{fixpoint attractor} of the orbit family.

\subsection{What D4 Achieves}

\begin{enumerate}[nosep]
\item \textbf{Observation Pipeline:} Every successful S7 emission feeds 
      artifact metadata into \texttt{NormRegistry::observe\_and\_learn()}.
\item \textbf{Complete Norm Synthesis:} Stub implementations in 
      \texttt{learning.rs} are replaced with operational code that produces 
      \texttt{NormLayers} with real artifacts.
\item \textbf{Norm-Guided Generation:} Matched norms inject structural 
      constraints into the HDAG before LLM calls, reducing token cost and 
      improving consistency.
\item \textbf{CLI Observability:} \texttt{isls norms} subcommand for 
      inspecting the norm catalog, candidates, and auto-discovered norms.
\end{enumerate}

\subsection{What D4 is NOT}

\begin{itemize}[nosep]
\item \textbf{Not iterative chat.} D4 is still single-shot generation. 
      Multi-turn refinement is D4+/D5.
\item \textbf{Not repository scraping.} External norm sources are D5.
\item \textbf{Not self-generation.} The system does not generate itself. 
      That is D6.
\item \textbf{Not a web UI.} CLI only. Studio integration comes later.
\end{itemize}


% ══════════════════════════════════════════════════════════════════════════════
\section{Architecture Overview}
% ══════════════════════════════════════════════════════════════════════════════

\subsection{Current State (D3)}

\begin{lstlisting}[style=bash]
Chat/TOML --> isls-chat --> isls-forge-llm (S0-S7) --> Output
                                                        |
                                                     (discarded)
\end{lstlisting}

The generation pipeline is open-loop. Generated artifacts vanish after 
emission. The \texttt{isls-norms} crate contains data structures and 
stubs for self-learning, but:

\begin{itemize}[nosep]
\item \texttt{build\_norm\_layers()} returns \texttt{NormLayers::default()} 
      --- empty.
\item \texttt{infer\_parameters()} returns \texttt{vec![]} --- no parameters.
\item \texttt{generalize\_name()} uses hardcoded string replacements for 
      3 words only.
\item \texttt{observe\_and\_learn()} is never called from the pipeline.
\item \texttt{ForgePlan::from\_composed\_plan()} exists but is never invoked.
\end{itemize}

\subsection{Target State (D4)}

\begin{lstlisting}[style=bash]
Chat/TOML --> isls-chat ---------> isls-forge-llm (S0-S7) --> Output
                |                       |      ^                 |
                |   norm-match          |      |                 |
                +-- NormRegistry <------+------+----- observe ---+
                        |                      |
                        +-- ~/.isls/norms.json -+
\end{lstlisting}

The generation pipeline becomes a closed loop:

\begin{enumerate}[nosep]
\item \textbf{Pre-generation:} \texttt{isls-chat} matches the user 
      description against the \texttt{NormRegistry}. Matched norms produce 
      a \texttt{ComposedPlan} that enriches the \texttt{AppSpec} with 
      norm-derived fields, validations, and business rules.
\item \textbf{Generation:} The HDAG pipeline runs as before (S0--S7), 
      but the \texttt{AppSpec} is now richer due to norm contributions.
\item \textbf{Post-generation:} After S7 Emit, artifact metadata is 
      collected from the generated files and fed into 
      \texttt{observe\_and\_learn()}.
\item \textbf{Persistence:} Candidates and auto-norms persist to 
      \texttt{\textasciitilde/.isls/norms.json}.
\end{enumerate}


% ══════════════════════════════════════════════════════════════════════════════
\section{Work Packages}
% ══════════════════════════════════════════════════════════════════════════════

D4 consists of four work packages executed sequentially. Each depends on 
the previous.

% ──────────────────────────────────────────────────────────────────────────────
\subsection{W1: Observation Pipeline}
\label{sec:w1}
% ──────────────────────────────────────────────────────────────────────────────

\subsubsection{Goal}

After every successful S7 emission, collect artifact metadata from the 
generated file tree and feed it into \texttt{NormRegistry::observe\_and\_learn()}.

\subsubsection{New File}

\begin{longtable}{lp{10cm}}
\toprule
\textbf{File} & \textbf{Purpose} \\
\midrule
\texttt{isls-forge-llm/src/artifact\_collector.rs} 
  & Walks the generated project directory after S7 Emit, inspects each 
    generated Rust file, and produces a \texttt{Vec<ObservedArtifact>}. \\
\bottomrule
\end{longtable}

\subsubsection{Artifact Collection Algorithm}

\begin{enumerate}[nosep]
\item Walk the \texttt{backend/src/} directory recursively.
\item For each \texttt{.rs} file, determine its \texttt{LayerType} from 
      the path:
      \begin{itemize}[nosep]
      \item \texttt{models/*.rs} $\to$ \texttt{LayerType::Model}
      \item \texttt{database/*\_queries.rs} $\to$ \texttt{LayerType::Query}
      \item \texttt{services/*.rs} $\to$ \texttt{LayerType::Service}
      \item \texttt{api/*.rs} $\to$ \texttt{LayerType::Api}
      \end{itemize}
\item For each file, extract:
      \begin{itemize}[nosep]
      \item \textbf{artifact\_type:} \texttt{"struct"} for models, 
            \texttt{"fn"} for queries/services/api.
      \item \textbf{name:} The struct name or function names found in 
            the file (regex: \texttt{pub struct (\textbackslash w+)} and 
            \texttt{pub (async )?fn (\textbackslash w+)}).
      \item \textbf{signature:} SHA-256 of the file content.
      \item \textbf{field\_names:} For structs, extract field names 
            (regex on struct body). For functions, extract parameter names.
      \end{itemize}
\item Also inspect:
      \begin{itemize}[nosep]
      \item \texttt{migrations/*.sql} $\to$ \texttt{LayerType::Database}, 
            extract table names from \texttt{CREATE TABLE} statements.
      \item \texttt{frontend/pages/*.js} $\to$ \texttt{LayerType::Frontend}, 
            extract page/component names.
      \item \texttt{api\_tests.rs} $\to$ \texttt{LayerType::Test}, 
            extract test function names.
      \end{itemize}
\item Return \texttt{Vec<ObservedArtifact>}.
\end{enumerate}

\subsubsection{Pipeline Integration}

In \texttt{staged\_closure.rs}, after the S7 Emit block, add the 
observation hook:

\begin{lstlisting}[style=rust]
// S7.1: Observe -- feed artifacts into norm learning
let collector = ArtifactCollector::new(&self.output_dir);
let observed = collector.collect()?;
let domain = &spec.app_name;
let run_id = format!("{}_{}", domain,
    chrono::Utc::now().format("%Y%m%d_%H%M%S"));
let mut registry = NormRegistry::new();
registry.observe_and_learn(&observed, domain, &run_id);
// registry.save() is called internally by observe_and_learn
\end{lstlisting}

\textbf{Constraint:} The observation MUST NOT block or delay S7 Emit. 
If collection fails, log a warning and continue --- graceful degradation 
(Rule~10).

\subsubsection{Modified Files}

\begin{longtable}{lp{10cm}}
\toprule
\textbf{File} & \textbf{Change} \\
\midrule
\texttt{isls-forge-llm/src/staged\_closure.rs}
  & Add observation hook after S7 Emit. Add \texttt{mod artifact\_collector} 
    to imports. \\
\texttt{isls-forge-llm/src/lib.rs}
  & Add \texttt{pub mod artifact\_collector;} \\
\texttt{isls-forge-llm/Cargo.toml}
  & Add dependency: \texttt{chrono = \{version = "0.4", 
    features = ["clock"]\}}, \texttt{sha2 = "0.10"} 
    (sha2 may already be present via isls-norms). \\
\bottomrule
\end{longtable}

\subsubsection{W1 Verification}

\begin{enumerate}[nosep]
\item Run \texttt{isls forge-chat -m "Pet shop with animals, owners, 
      appointments" --api-key \$KEY --output ./output/petshop}
\item Verify \texttt{\textasciitilde/.isls/norms.json} exists and contains 
      at least one candidate entry with \texttt{status: "Observing"}.
\item Verify the candidate has \texttt{domains: ["petshop"]} and 
      \texttt{observation\_count: 1}.
\item Run a second domain (e.g. hotel). Verify the same candidate now has 
      2 domains if patterns overlap.
\end{enumerate}


% ──────────────────────────────────────────────────────────────────────────────
\subsection{W2: Complete Norm Synthesis}
\label{sec:w2}
% ──────────────────────────────────────────────────────────────────────────────

\subsubsection{Goal}

Replace the three stub implementations in \texttt{learning.rs} so that 
promoted candidates produce norms with \emph{real} \texttt{NormLayers} 
containing artifact templates usable by the composition engine.

\subsubsection{Stubs to Complete}

\paragraph{1. \texttt{build\_norm\_layers()}}

Current: returns \texttt{NormLayers::default()} (all fields empty).

Target: Iterate over \texttt{candidate.common\_artifacts}. For each 
\texttt{AbstractedArtifact}, create the corresponding layer artifact:

\begin{itemize}[nosep]
\item \texttt{LayerType::Database} $\to$ \texttt{DatabaseArtifact} with 
      \texttt{table = name\_template}, columns from \texttt{common\_fields} 
      mapped to SQL types via a simple type-inference heuristic 
      (field name contains ``id'' $\to$ \texttt{INTEGER}, 
      ``name''/``title''/``description'' $\to$ \texttt{TEXT}, 
      ``price''/``amount''/``quantity'' $\to$ \texttt{NUMERIC}, 
      ``date''/``time''/``created''/``updated'' $\to$ \texttt{TIMESTAMPTZ}, 
      ``active''/``is\_'' $\to$ \texttt{BOOLEAN}, 
      default $\to$ \texttt{TEXT}).
\item \texttt{LayerType::Model} $\to$ \texttt{ModelArtifact} with 
      \texttt{struct\_name = name\_template} (PascalCase), fields from 
      \texttt{common\_fields} with Rust type inference 
      (same heuristic as above but mapped to Rust types: \texttt{i32}, 
      \texttt{String}, \texttt{f64}, \texttt{chrono::NaiveDateTime}, 
      \texttt{bool}).
\item \texttt{LayerType::Query} $\to$ \texttt{QueryArtifact} with 
      standard CRUD function names: 
      \texttt{get\_\{entity\}}, \texttt{list\_\{entities\}}, 
      \texttt{create\_\{entity\}}, \texttt{update\_\{entity\}}, 
      \texttt{delete\_\{entity\}} (per Rule~3).
\item \texttt{LayerType::Service} $\to$ \texttt{ServiceArtifact} with 
      thin wrapper names matching query functions.
\item \texttt{LayerType::Api} $\to$ \texttt{ApiArtifact} with scope 
      \texttt{/api/\{entities\}} and standard CRUD handlers.
\item \texttt{LayerType::Frontend} $\to$ \texttt{FrontendArtifact} with 
      a page template name.
\item \texttt{LayerType::Test} $\to$ \texttt{TestArtifact} with CRUD 
      test names.
\end{itemize}

Artifacts with \texttt{confidence < min\_artifact\_presence} (from 
\texttt{PromotionCriteria}) are excluded.

\paragraph{2. \texttt{infer\_parameters()}}

Current: returns \texttt{vec![]}.

Target: Scan \texttt{common\_fields} across all artifacts. Fields 
whose \texttt{presence\_rate < 1.0} become optional 
\texttt{NormParameter}s (type = \texttt{ParamType::Bool}, name = 
\texttt{"include\_\{field\_name\}"}). Fields with \texttt{presence\_rate 
== 1.0} are fixed (not parameters).

\paragraph{3. \texttt{generalize\_name()}}

Current: hardcoded \texttt{.replace("product", "\{entity\}")} etc.

Target: Extract the entity-specific word by comparing the artifact name 
against the \texttt{observed\_on} field of the pattern. Replace any 
occurrence of the entity name (case-insensitive) with 
\texttt{"\{entity\}"}. Example: \texttt{"ProductModel"} observed on 
entity \texttt{"product"} $\to$ \texttt{"\{entity\}Model"}.

The replacement must also handle plurals: if \texttt{observed\_on} is 
\texttt{"product"}, also replace \texttt{"products"} with 
\texttt{"\{entities\}"}.

\subsubsection{Modified Files}

\begin{longtable}{lp{10cm}}
\toprule
\textbf{File} & \textbf{Change} \\
\midrule
\texttt{isls-norms/src/learning.rs}
  & Replace \texttt{build\_norm\_layers()}, \texttt{infer\_parameters()}, 
    \texttt{generalize\_name()} with operational implementations. \\
\bottomrule
\end{longtable}

\subsubsection{W2 Verification}

\begin{enumerate}[nosep]
\item Unit test: create 5+ \texttt{ObservedArtifact}s from 3 domains 
      with consistent structure. Feed into a \texttt{NormRegistry} 
      5 times (meeting \texttt{PromotionCriteria::default()}).
\item Assert a norm with ID prefix \texttt{ISLS-NORM-AUTO-} was created.
\item Assert the norm's \texttt{layers.database} is non-empty.
\item Assert the norm's \texttt{layers.model} is non-empty with correct 
      \texttt{struct\_name} containing \texttt{"\{entity\}"}.
\item Assert the norm's \texttt{layers.query} contains 5 CRUD function 
      templates.
\item Assert the norm has at least one \texttt{NormParameter} if any 
      field had \texttt{presence\_rate < 1.0}.
\end{enumerate}


% ──────────────────────────────────────────────────────────────────────────────
\subsection{W3: Norm-Guided Generation}
\label{sec:w3}
% ──────────────────────────────────────────────────────────────────────────────

\subsubsection{Goal}

When generating a new application, match the user's description against 
the \texttt{NormRegistry} (including auto-discovered norms) and use 
matched norms to enrich the \texttt{AppSpec} before HDAG construction.

\subsubsection{Integration Point: \texttt{isls-chat}}

Currently, \texttt{forge-chat} does:

\begin{lstlisting}[style=bash]
message --> LLM entity extraction --> JSON --> TOML --> forge-v2
\end{lstlisting}

After D4:

\begin{lstlisting}[style=bash]
message --> NormRegistry::match_description(message)
        --> ActivatedNorms (confidence >= 0.3)
        --> compose_norms_with_registry(activated, params, registry)
        --> ComposedPlan
        --> LLM entity extraction (enriched prompt with norm hints)
        --> JSON --> TOML (merge ComposedPlan fields into entities)
        --> forge-v2
\end{lstlisting}

\subsubsection{Norm Enrichment Algorithm}

\begin{enumerate}[nosep]
\item Load \texttt{NormRegistry} (includes builtin + auto-discovered).
\item Call \texttt{registry.match\_description(\&message)} $\to$ 
      \texttt{Vec<ActivatedNorm>}.
\item If any norms activated with confidence $\geq 0.5$:
      \begin{enumerate}[nosep,label=\alph*.]
      \item Compose: \texttt{compose\_norms\_with\_registry(\&activated, 
            \&params, \&registry)} $\to$ \texttt{ComposedPlan}.
      \item For each \texttt{ModelArtifact} in the plan: add suggested 
            fields to the LLM extraction prompt as hints 
            (``Common fields for \{entity\_type\}: ...'').
      \item After LLM extraction: merge any ComposedPlan fields that 
            the LLM did not already produce (additive only, never 
            override LLM output).
      \item Log activated norms: 
            \texttt{tracing::info!("D4 norm enrichment: \{norm\_ids\}")}.
      \end{enumerate}
\item If no norms activated: proceed as D3 (pure LLM extraction).
\end{enumerate}

\textbf{Critical constraint:} Norm enrichment is \emph{additive only}. 
It MUST NOT remove or override entities/fields produced by the LLM. 
The LLM has domain-specific context that norms lack. Norms provide 
\emph{structural priors}, not \emph{corrections}.

\subsubsection{New File}

\begin{longtable}{lp{10cm}}
\toprule
\textbf{File} & \textbf{Purpose} \\
\midrule
\texttt{isls-chat/src/norm\_enrichment.rs}
  & Contains \texttt{enrich\_with\_norms(message, entities) -> entities}. 
    Loads registry, matches, composes, and merges norm-derived fields 
    into the entity list. Returns enriched entities. \\
\bottomrule
\end{longtable}

\subsubsection{Modified Files}

\begin{longtable}{lp{10cm}}
\toprule
\textbf{File} & \textbf{Change} \\
\midrule
\texttt{isls-chat/src/lib.rs}
  & Add \texttt{pub mod norm\_enrichment;} \\
\texttt{isls-chat/src/extract.rs} (or equivalent)
  & After LLM entity extraction, call 
    \texttt{norm\_enrichment::enrich\_with\_norms()} before converting 
    to TOML. \\
\texttt{isls-chat/Cargo.toml}
  & Add dependency: \texttt{isls-norms = \{path = "../isls-norms"\}} \\
\bottomrule
\end{longtable}

\subsubsection{Prompt Enhancement}

When norms are activated, append to the LLM entity-extraction prompt:

\begin{lstlisting}[style=bash]
## Norm Hints (optional structural guidance)
The following patterns have been observed across similar applications.
Consider including these fields if they fit the domain:

- Entities with status tracking commonly include:
  status (String), created_at (DateTime), updated_at (DateTime)
- Entities with inventory commonly include:
  quantity (i32), unit_price (f64), reorder_level (i32)

These are suggestions, not requirements. Extract entities that fit
the user's description first, then add norm-suggested fields only
if they are relevant.
\end{lstlisting}

The exact hint text is generated from \texttt{ComposedPlan.models} 
at runtime --- not hardcoded.

\subsubsection{W3 Verification}

\begin{enumerate}[nosep]
\item Generate 3 domains (petshop, hotel, clinic) to accumulate 
      observations past the promotion threshold.
\item Generate a 4th domain (e.g. ``library management system'').
\item Verify that the generated app contains norm-suggested fields 
      that were NOT in the user's one-sentence description.
\item Verify the norm enrichment is logged: 
      \texttt{grep "D4 norm enrichment" <output>}.
\item Verify compilation still passes (\texttt{cargo build} succeeds).
\end{enumerate}


% ──────────────────────────────────────────────────────────────────────────────
\subsection{W4: CLI Observability}
\label{sec:w4}
% ──────────────────────────────────────────────────────────────────────────────

\subsubsection{Goal}

Add an \texttt{isls norms} subcommand for inspecting the norm ecosystem 
without modifying generation behavior.

\subsubsection{Commands}

\begin{lstlisting}[style=bash]
isls norms list              # List all norms (builtin + auto)
isls norms list --auto-only  # Only auto-discovered norms
isls norms inspect <norm-id> # Show full norm details
isls norms candidates        # List candidate pool
isls norms stats             # Summary statistics
isls norms reset             # Delete ~/.isls/norms.json (with confirm)
\end{lstlisting}

\paragraph{Output Format (example for \texttt{isls norms list}):}

\begin{lstlisting}[style=bash]
ISLS Norm Catalog (27 builtin, 2 auto-discovered)
---
ISLS-NORM-0042  CRUD-Entity       builtin   Molecule  [Db,Model,Query,Svc,Api]
ISLS-NORM-0088  JWT-Auth          builtin   Molecule  [Db,Model,Query,Svc,Api]
...
ISLS-NORM-AUTO-0001  Auto-StateMachine  auto  Molecule  [Db,Model,Query,Svc]
    domains: warehouse, ecommerce, projectmgmt
    observations: 7, consistency: 0.91
ISLS-NORM-AUTO-0002  Auto-Inventory     auto  Molecule  [Db,Model,Query]
    domains: warehouse, petshop, restaurant
    observations: 5, consistency: 0.87
\end{lstlisting}

\paragraph{Output Format (example for \texttt{isls norms stats}):}

\begin{lstlisting}[style=bash]
ISLS Norm Statistics
---
Builtin norms:       24 (20 molecule, 3 organism, 1 meta)
Auto-discovered:     2
Candidates:          5 (3 observing, 2 promoted)
Total observations:  23 across 6 domains
Wirings:             8
Persistence:         ~/.isls/norms.json (4.2 KB)
\end{lstlisting}

\subsubsection{Modified Files}

\begin{longtable}{lp{10cm}}
\toprule
\textbf{File} & \textbf{Change} \\
\midrule
\texttt{isls-cli/src/main.rs}
  & Add \texttt{Command::Norms} variant with subcommands. Parse 
    \texttt{isls norms <subcmd>} from args. \\
\texttt{isls-cli/src/cmd\_norms.rs} (new)
  & Implements \texttt{cmd\_norms\_list()}, \texttt{cmd\_norms\_inspect()}, 
    \texttt{cmd\_norms\_candidates()}, \texttt{cmd\_norms\_stats()}, 
    \texttt{cmd\_norms\_reset()}. \\
\texttt{isls-cli/Cargo.toml}
  & Add dependency: \texttt{isls-norms = \{path = "../isls-norms"\}} \\
\bottomrule
\end{longtable}

\subsubsection{W4 Verification}

\begin{enumerate}[nosep]
\item \texttt{isls norms list} --- shows builtin norms, compiles, 
      no panic.
\item \texttt{isls norms stats} --- shows summary, numbers match 
      actual registry state.
\item \texttt{isls norms candidates} --- shows candidate pool after 
      at least one generation run.
\item \texttt{isls norms inspect ISLS-NORM-0042} --- shows full 
      norm details including layers and triggers.
\end{enumerate}


% ══════════════════════════════════════════════════════════════════════════════
\section{File Inventory}
% ══════════════════════════════════════════════════════════════════════════════

\subsection{New Files}

\begin{longtable}{llp{7cm}}
\toprule
\textbf{Crate} & \textbf{File} & \textbf{Purpose} \\
\midrule
isls-forge-llm & \texttt{src/artifact\_collector.rs} 
  & Post-emission artifact walker \\
isls-chat & \texttt{src/norm\_enrichment.rs} 
  & Pre-generation norm matching + enrichment \\
isls-cli & \texttt{src/cmd\_norms.rs} 
  & \texttt{isls norms} subcommand handlers \\
\bottomrule
\end{longtable}

\subsection{Modified Files}

\begin{longtable}{llp{7cm}}
\toprule
\textbf{Crate} & \textbf{File} & \textbf{Change} \\
\midrule
isls-forge-llm & \texttt{src/staged\_closure.rs} 
  & Add S7.1 observation hook \\
isls-forge-llm & \texttt{src/lib.rs} 
  & Add \texttt{pub mod artifact\_collector} \\
isls-forge-llm & \texttt{Cargo.toml} 
  & Add chrono dependency \\
isls-norms & \texttt{src/learning.rs} 
  & Complete 3 stub functions \\
isls-chat & \texttt{src/lib.rs} 
  & Add \texttt{pub mod norm\_enrichment} \\
isls-chat & \texttt{src/extract.rs} 
  & Call enrichment after LLM extraction \\
isls-chat & \texttt{Cargo.toml} 
  & Add isls-norms dependency \\
isls-cli & \texttt{src/main.rs} 
  & Add Norms command variant \\
isls-cli & \texttt{Cargo.toml} 
  & Add isls-norms dependency \\
\bottomrule
\end{longtable}

\subsection{Unchanged Files}

All 48 structural file generators, the HDAG builder, the Oracle trait, 
S6 Coagula, and the full isls-norms type system (\texttt{types.rs}, 
\texttt{catalog.rs}, \texttt{composition.rs}) remain \textbf{unchanged}. 
D4 extends the system without modifying the proven generation core.


% ══════════════════════════════════════════════════════════════════════════════
\section{The 12 Rules --- D4 Compliance}
% ══════════════════════════════════════════════════════════════════════════════

\begin{longtable}{clp{9cm}}
\toprule
\textbf{\#} & \textbf{Status} & \textbf{D4 Impact} \\
\midrule
1  & \color{isls-green}\checkmark & Norm enrichment adds fields to EntityDef. 
     Structural files remain structural. LLM surface unchanged (7 files). \\
2  & \color{isls-green}\checkmark & No new modules violate re-export rule. \\
3  & \color{isls-green}\checkmark & \texttt{build\_norm\_layers} generates 
     CRUD names per Rule~3 convention. \\
4  & \color{isls-green}\checkmark & No change to auth handling. \\
5  & \color{isls-green}\checkmark & No change to query signatures. \\
6  & \color{isls-green}\checkmark & No change to query files. \\
7  & \color{isls-green}\checkmark & No new crypto or logging dependencies. \\
8  & \color{isls-green}\checkmark & Norm-enriched fields become part of 
     EntityDef, so seed INSERT includes them. \\
9  & \color{isls-green}\checkmark & No Dockerfile change. \\
10 & \color{isls-green}\checkmark & Artifact collection failure $\to$ warn 
     + continue. Norm matching failure $\to$ proceed as D3. \\
11 & \color{isls-green}\checkmark & ProvidedSymbol = structural content. 
     Unchanged. \\
12 & \color{isls-green}\checkmark & LLM prompts get norm hints as 
     additional context --- more type context, not less. \\
\bottomrule
\end{longtable}


% ══════════════════════════════════════════════════════════════════════════════
\section{Data Flow: Complete D4 Cycle}
% ══════════════════════════════════════════════════════════════════════════════

\begin{enumerate}
\item \textbf{User:} \texttt{isls forge-chat -m "Hotel management 
      with rooms, guests, bookings, and housekeeping" --api-key \$KEY}

\item \textbf{isls-chat (Norm Match):}
      \begin{enumerate}[nosep,label=\alph*.]
      \item Load \texttt{NormRegistry} from \texttt{\textasciitilde/.isls/norms.json}.
      \item \texttt{match\_description("Hotel management...")} $\to$ 
            activated norms (e.g. ISLS-NORM-0042 CRUD-Entity at 0.75, 
            ISLS-NORM-AUTO-0001 Auto-StateMachine at 0.50).
      \item \texttt{compose\_norms\_with\_registry()} $\to$ ComposedPlan 
            with suggested fields (status, created\_at, updated\_at).
      \end{enumerate}

\item \textbf{isls-chat (LLM Extraction):}
      \begin{enumerate}[nosep,label=\alph*.]
      \item Build extraction prompt with norm hints appended.
      \item LLM returns JSON entities: Room, Guest, Booking, 
            HousekeepingTask.
      \item \texttt{enrich\_with\_norms()} merges norm-suggested fields 
            (e.g. adds \texttt{status} to HousekeepingTask if not already 
            present).
      \end{enumerate}

\item \textbf{isls-chat (TOML):} Convert enriched entities to TOML.

\item \textbf{isls-forge-llm (S0--S7):} Standard HDAG pipeline with 
      enriched AppSpec. Generated app has norm-informed fields.

\item \textbf{isls-forge-llm (S7.1 Observe):}
      \begin{enumerate}[nosep,label=\alph*.]
      \item \texttt{ArtifactCollector} walks generated project.
      \item Produces \texttt{Vec<ObservedArtifact>} (e.g. 28 artifacts 
            across 6 layers).
      \item \texttt{registry.observe\_and\_learn(\&artifacts, "hotel", 
            \&run\_id)}.
      \item Cross-layer patterns extracted, matched against candidates.
      \item Candidates updated (observation\_count++, domains extended).
      \item If any candidate meets PromotionCriteria $\to$ auto-promoted.
      \item Persisted to disk.
      \end{enumerate}

\item \textbf{User verifies:} \texttt{cd output/hotel/backend \&\& 
      cargo build} --- compiles. Norm-enriched fields visible in 
      generated models.
\end{enumerate}


% ══════════════════════════════════════════════════════════════════════════════
\section{Promotion Criteria}
% ══════════════════════════════════════════════════════════════════════════════

The existing \texttt{PromotionCriteria::default()} values are retained:

\begin{longtable}{lrl}
\toprule
\textbf{Criterion} & \textbf{Threshold} & \textbf{Meaning} \\
\midrule
\texttt{min\_consistency}        & 0.85  & Jaccard similarity of layer 
                                            sets across observations \\
\texttt{min\_domains}            & 3     & Observed in $\geq$3 distinct 
                                            app domains \\
\texttt{min\_layers}             & 4     & Pattern spans $\geq$4 
                                            layers consistently \\
\texttt{min\_observations}       & 5     & At least 5 total 
                                            observations \\
\texttt{min\_artifact\_presence} & 0.80  & Each artifact template 
                                            present in $\geq$80\% of 
                                            observations \\
\bottomrule
\end{longtable}

These thresholds prevent premature promotion. With 6 verified apps 
(Warehouse, E-Commerce, Projektmanagement, Restaurant, Library, Gym) 
already generated, and D4 adding observation, the first auto-norms 
should emerge after $\sim$2--3 additional generation runs (total 
$\geq$5 observations across $\geq$3 domains).


% ══════════════════════════════════════════════════════════════════════════════
\section{Error Handling}
% ══════════════════════════════════════════════════════════════════════════════

\begin{longtable}{lp{10cm}}
\toprule
\textbf{Failure Mode} & \textbf{Behavior} \\
\midrule
Artifact collection IO error & Log warning, skip observation. 
  App is still emitted successfully. \\
\texttt{norms.json} corrupt & Log warning, start with fresh 
  registry (builtin norms only). \\
Norm matching returns empty & Proceed as D3 (pure LLM extraction). \\
Composition conflict & Log warning, fall back to D3 path. \\
Enrichment produces invalid field types & Validation in 
  \texttt{json\_to\_toml()} catches it. Log + skip invalid fields. \\
Auto-norm has empty NormLayers & Skip during enrichment 
  (confidence filter). \\
\bottomrule
\end{longtable}

All error paths converge on: \textbf{generate the app anyway, as D3 
would.} The norm system is strictly additive --- it can enhance but 
never block generation.


% ══════════════════════════════════════════════════════════════════════════════
\section{Test Suite}
% ══════════════════════════════════════════════════════════════════════════════

\subsection{Unit Tests (in-crate)}

\begin{longtable}{lp{10cm}}
\toprule
\textbf{Test} & \textbf{Assertion} \\
\midrule
\texttt{test\_artifact\_collector\_models} 
  & Collector extracts struct names + fields from a mock models/ dir. \\
\texttt{test\_artifact\_collector\_queries} 
  & Collector extracts function names from mock query files. \\
\texttt{test\_artifact\_collector\_layers} 
  & Correct LayerType assignment for each directory. \\
\texttt{test\_build\_norm\_layers\_database} 
  & \texttt{build\_norm\_layers} produces non-empty DatabaseArtifact 
    with correct column types. \\
\texttt{test\_build\_norm\_layers\_model} 
  & Produces ModelArtifact with \texttt{\{entity\}} in struct\_name. \\
\texttt{test\_build\_norm\_layers\_crud} 
  & QueryArtifact has 5 CRUD functions per Rule~3. \\
\texttt{test\_infer\_parameters} 
  & Fields with presence\_rate $<$ 1.0 become NormParameters. \\
\texttt{test\_generalize\_name\_dynamic} 
  & Entity-specific words replaced with \texttt{\{entity\}}, 
    plurals with \texttt{\{entities\}}. \\
\texttt{test\_full\_promotion\_cycle} 
  & 5 observations across 3 domains $\to$ candidate promoted $\to$ 
    auto-norm has non-empty layers. \\
\texttt{test\_norm\_enrichment\_additive} 
  & Enrichment adds fields but never removes LLM-produced ones. \\
\texttt{test\_norm\_enrichment\_fallback} 
  & If no norms match, entities pass through unchanged. \\
\bottomrule
\end{longtable}

\subsection{Integration Tests}

\begin{longtable}{lp{10cm}}
\toprule
\textbf{Test} & \textbf{Assertion} \\
\midrule
\texttt{test\_d4\_observation\_e2e} 
  & Run \texttt{forge-chat} with mock oracle. Verify 
    \texttt{norms.json} updated. \\
\texttt{test\_d4\_norms\_cli} 
  & Run \texttt{isls norms list} and \texttt{isls norms stats}. 
    Verify output format, no panic. \\
\bottomrule
\end{longtable}

\subsection{Acceptance Test (Manual)}

\begin{lstlisting}[style=bash]
# Step 1: Generate 3+ domains to accumulate observations
isls forge-chat -m "Pet shop with animals, owners, appointments" \
    --api-key $KEY --output ./output/petshop
isls forge-chat -m "Hotel with rooms, guests, bookings" \
    --api-key $KEY --output ./output/hotel
isls forge-chat -m "Clinic with patients, doctors, appointments" \
    --api-key $KEY --output ./output/clinic

# Step 2: Verify candidates exist
isls norms candidates
# Expected: >= 1 candidate with status Observing, domains >= 3

# Step 3: Generate 2 more domains to trigger promotion
isls forge-chat -m "Gym with members, trainers, classes" \
    --api-key $KEY --output ./output/gym2
isls forge-chat -m "School with students, teachers, courses" \
    --api-key $KEY --output ./output/school

# Step 4: Verify auto-norm emergence
isls norms list --auto-only
# Expected: >= 1 auto-discovered norm (ISLS-NORM-AUTO-*)

# Step 5: Generate with norm enrichment
isls forge-chat -m "Spa with treatments, therapists, bookings" \
    --api-key $KEY --output ./output/spa
cd output/spa/backend && cargo build && cd ../../..
# Expected: compiles, generated models contain norm-enriched fields

# Step 6: Verify enrichment was applied
grep -r "status" output/spa/backend/src/models/
# Expected: status fields in entities where Auto-StateMachine applied
\end{lstlisting}


% ══════════════════════════════════════════════════════════════════════════════
\section{Constraints for Claude Code}
% ══════════════════════════════════════════════════════════════════════════════

\begin{enumerate}
\item \textbf{Do not modify} \texttt{structural.rs}, \texttt{hdag.rs}, 
      \texttt{oracle.rs}, or any of the 48 structural file generators.
\item \textbf{Do not modify} \texttt{types.rs} or \texttt{catalog.rs} 
      in \texttt{isls-norms} --- the type system is stable.
\item \textbf{composition.rs} is read-only --- the composition engine 
      is complete.
\item The \texttt{artifact\_collector.rs} MUST use only \texttt{std::fs} 
      and regex --- no external parsing crates, no tokio.
\item \texttt{norm\_enrichment.rs} MUST be pure (no IO except loading 
      the registry). Input: message + entities. Output: enriched entities.
\item The observation hook MUST NOT delay or block S7 Emit. Wrap in 
      a catch-all that logs warnings on failure.
\item All new code MUST compile with \texttt{cargo build} on the 
      existing workspace without new \texttt{[patch]} or 
      \texttt{[workspace.dependencies]} changes (except adding 
      chrono and path dependencies as specified).
\item Test with \texttt{cargo test -p isls-norms} and 
      \texttt{cargo test -p isls-forge-llm} --- both must pass.
\item Branch: \texttt{claude/implement-isls-d4-spec}.
\item Commit message format: \texttt{D4/W\{n\}: <description>} 
      (e.g. \texttt{D4/W1: Add artifact collector and observation hook}).
\end{enumerate}


% ══════════════════════════════════════════════════════════════════════════════
\section{Dependency Graph}
% ══════════════════════════════════════════════════════════════════════════════

\begin{verbatim}
W1 (Observation Pipeline)
  |
  v
W2 (Complete Norm Synthesis)
  |
  v
W3 (Norm-Guided Generation)
  |
  v
W4 (CLI Observability)
\end{verbatim}

W1 must work before W2 can be tested end-to-end. W2 must produce 
real norms before W3 can use them. W4 is independent but most useful 
after W1--W3 populate the registry.

Recommended implementation order: W1 $\to$ W2 $\to$ W4 $\to$ W3. 
This allows testing the observation + synthesis loop (W1+W2) before 
wiring norm-guided generation (W3), with W4 providing observability 
tools during development.


% ══════════════════════════════════════════════════════════════════════════════
\section{Token Budget Estimate}
% ══════════════════════════════════════════════════════════════════════════════

D4 adds \textbf{zero additional LLM tokens per generation run} in the 
default path (no norms matched). When norms are matched, the prompt 
enhancement adds $\sim$200--500 tokens to the entity extraction call 
(negligible at \$0.15/run).

The observation step (W1) is pure Rust file IO --- zero tokens.

Total API cost impact: $<$\$0.02/run additional.


% ══════════════════════════════════════════════════════════════════════════════
\section{Post-D4 Outlook}
% ══════════════════════════════════════════════════════════════════════════════

With D4 complete, the norm system becomes self-reinforcing:

\begin{itemize}[nosep]
\item Each generation run is an observation.
\item Observations accumulate into candidates.
\item Candidates promote to norms.
\item Norms enrich future generations.
\item Enriched generations produce more consistent artifacts.
\item More consistent artifacts promote faster.
\end{itemize}

This is the \emph{Ophan orbit} closing into a fixpoint: the system's 
generative topology converges toward its own attractors.

D5 (Repository Scraping) extends the observation surface beyond 
self-generated apps to external repositories. D6 (M\"obius-Umst\"ulpung) 
uses the mature norm catalog to generate ISLS itself.


% ══════════════════════════════════════════════════════════════════════════════
\section{Verification Sequence (Complete)}
\label{sec:verification}
% ══════════════════════════════════════════════════════════════════════════════

\begin{enumerate}
\item \texttt{cargo build} --- entire workspace compiles.
\item \texttt{cargo test -p isls-norms} --- all unit tests pass 
      (including new W2 tests).
\item \texttt{cargo test -p isls-forge-llm} --- all unit tests pass 
      (including W1 collector tests).
\item \texttt{cargo test -p isls-chat} --- enrichment tests pass.
\item \texttt{cargo test -p isls-cli} --- norms CLI tests pass.
\item \texttt{isls norms list} --- outputs norm catalog without panic.
\item \texttt{isls norms stats} --- outputs statistics.
\item \texttt{isls forge-chat -m "..." --api-key \$KEY --output ./test} 
      --- generates app.
\item \texttt{cd test/backend \&\& cargo build} --- generated app compiles.
\item \texttt{\textasciitilde/.isls/norms.json} --- contains candidate entries.
\item After 5+ runs across 3+ domains: \texttt{isls norms list --auto-only} 
      shows at least one auto-discovered norm.
\end{enumerate}

\vfill
\begin{center}
\color{isls-gray}\rule{0.5\textwidth}{0.4pt}\\[0.5em]
{\small End of D4 Specification.}
\end{center}

\end{document}
